\documentclass[small]{cv}

% 국문판(32mm)과 맞춤. 기본 24mm로는 2025.03--present가 두 줄로 접혀
% 항목마다 한 줄씩 낭비되고, 국문판과 좌측 배치도 달라 보임.
\setdatewidth{30mm}

% [확인] 아래 표시된 항목은 본인 공개 페이지(GitHub 프로필, til.99jik.com, selab.knu.ac.kr)에서
% 끌어온 추론. 사실 확인 후 주석 지울 것.

\cvname{Jeongin Kim}
% 자율주행은 뺐음. 읽은 논문에서 추론한 것이고 실적이 없음.
% 실제 실적은 SW 테스팅(KCC 25/26), SIL(ABB), 산업 공정 데이터 ML(UCWIT 25).
\cvtagline{M.S. student in software testing. LLM-driven generation of test environments,
and predictive modeling on industrial process data.}
% 전화번호는 뺐음. 저장소가 public이라 push하면 그대로 노출됨.
% 제출할 때 필요하면 그때 한 줄 넣고 빌드할 것.
% 링크는 둘만. GitHub은 프로젝트 코드 검증 경로라 필수, 개인 사이트는 상위 도메인 하나로 충분.
% 거주지도 뺐음. 학력과 경력 항목마다 Daegu, Korea가 이미 붙어 있어 중복이고,
% 학술 CV에서 거주지는 관례적으로 안 씀.
\cvcontact{%
  \cvmail{99jik@99jik.com}\cvsep
  \cvlink{https://99jik.com}{99jik.com}\cvsep
  \cvlink{https://github.com/99JIK}{github.com/99JIK}%
}

\begin{document}
\makecvheader

% ============================================================ Education
\section{Education}
% [확인] 지도교수/연구실은 GitHub bio + selab.knu.ac.kr에서 추론
\cventry{2025.03--present}{M.S. Student, Computer Science and Engineering}
  {Kyungpook National University}{Daegu, Korea}
  {Advisor: Prof. Woojin Lee, Software Testing Lab (STLAB). Thesis topic not yet fixed.}

% 성적증명서 기준. 생년월일, 학번, 야간 여부는 CV에 안 넣음.
% 학점은 지역 옆 작은 muted 슬롯으로. 두 학위 모두 표기해 비대칭을 없앰.
% 하나만 적고 하나를 비우면 심사자가 왜 하나만 없는지 알아챔. 숨기는 게 더 불리함.
% (Advanced Major Course)는 소속 줄로 내렸음. 제목에 두면 길어져서 오른쪽 GPA 표기와
% 충돌해 "GPA"만 첫 줄에 남고 나머지가 다음 줄로 흘러 소속 줄을 덮었음.
\cventry{2023.03--2024.02}{B.Eng., Computer Information Engineering}
  {Yeungjin University \cvnote{(Advanced Major Course)}}
  {GPA 3.50/4.50\cvsep Daegu, Korea}{}

% 계열은 소속 줄로 내렸음. 제목에 두면 길어져서 오른쪽 표기와 충돌해 줄이 깨짐.
\cventry{2018.03--2023.02}{Associate Degree in Engineering, Web Database}
  {Computer Information Division, Yeungjin University}{GPA 3.05/4.50\cvsep Daegu, Korea}{}

% 5년 재학의 이유가 여기서 바로 설명돼야 함. 안 밝히면 심사자가 먼저 묻는다.
% 날짜는 본인 기억 기준. 병적증명서(정부24)로 확인 필요.
\cvline{2018.10--2020.06}{Republic of Korea Army, mandatory military service (completed)}

% ============================================================ Research Interests
\section{Research Interests}
% 논문 세 편과 ABB 과제에서 실제로 다룬 것만. 자율주행과 퍼징은 실적이 없어 뺐음.
Software testing, LLM-based test generation, SIL verification environments,
predictive modeling on industrial process data

% ============================================================ Publications
\section{Publications \cvsectionnote{\corr corresponding author, own name in bold}}

% 최신순으로 위에서 아래로. 번호는 1부터 올라감. 서지사항은 DBpia 등재 기준.
% DBpia 목록에 교신저자가 안 뜨므로 이우진은 따로 붙였음.
% 영문 표기는 KCC 스타일로 고정: Jeongin Kim, Woojin Lee (붙여 쓰고 두 번째 음절 소문자).
% UCWIT 원문은 Jeong In Kim / Woo Jin Lee로 실렸으나 CV에서는 KCC 쪽으로 맞췄음.
% 표기가 갈리면 Scholar나 DBpia에서 같은 저자로 안 묶임. 앞으로 투고할 때도 이걸로 통일할 것.
\subsection*{Domestic Conference Papers \cvsectionnote{all in Korean}}
\begin{publications}
  \item \me{Jeongin Kim}, Deokyeop Kim, Hyeryung Suh, and Woojin Lee\corr.
        Automatic Generation of Test Environment Behaviors Using LLM for
        Embedded Software Testing.
        \venue{Proc. Korea Computer Congress (KCC)}, pp. 582--584, June 2026.
        \cvnote{Code: \cvlink{https://github.com/99JIK/KCC2026}{github.com/99JIK/KCC2026}}

  % 영문 제목은 논문 원문에 실린 공식 제목.
  \item \me{Jeongin Kim} and Woojin Lee\corr.
        A Study on the Utilization of Sampling Time-point Data for Performance
        Enhancement of the Biodegradable Fiber Spinning Process.
        \venue{Symposium on Ubiquitous Computing and Web Information Technology (UCWIT)},
        KIISE Yeongnam Branch, November 2025.

  % DBpia 원문에 공식 영문 제목 없음. 아래는 내가 옮긴 것.
  \item \me{Jeongin Kim} and Woojin Lee\corr.
        The Ability of Generative AI to Understand and Utilize Test Cases.
        \venue{Proc. Korea Computer Congress (KCC)}, pp. 397--399, July 2025.
\end{publications}

% ============================================================ Patents
% 가출원 상태를 반드시 명시할 것. 정식 출원 전이고 심사도 안 들어갔음.
% 기여도 수치(30/20/50)는 안 적음. 관례상 발명자 순서만 씀.
\section{Patents}
% 발명자 순서는 출원서 기준으로 본인이 확인해 준 것(김정인 먼저).
% 특허사무소 송부 문서의 대표 발명자란은 이우진이 먼저인데, 그 필드는 연락 책임자를
% 앞세우는 사무소 관행이라 출원서 순서와 다름. 문서 순서로 적었다가 되돌렸음.
% 정식 출원번호 나오면 추가할 것.
\begin{publications}
  \item \me{Jeongin Kim}, Deokyeop Kim, and Woojin Lee.
        A Method for Selecting Robust Optimal Process Parameter Combinations.
        \cvnote{Provisional filed, full application in progress (July 2026).
        Applicant: KNU Industry-Academic Cooperation Foundation. MOTIE/KIAT P0022335.}
\end{publications}

% ============================================================ Experience
\section{Experience}
% 석사 입학은 2025.03이지만 연구실 합류는 2024.12. 사실대로 적으면 공백도 6개월로 줄어듦.
\cventry{2024.12--present}{Graduate Researcher}
  {Software Testing Lab, Kyungpook National University}{Daegu, Korea}
  {\begin{cvitems}
     \item LLM-driven generation of environment object behaviors for SIL testing. A
           three-stage decomposition-induction-synthesis prompt pipeline more than doubled
           coverage over zero-shot and few-shot baselines across 11 objects in three
           domains (elevator, drone, smart factory), consistent on GPT-4 and
           Llama-3 70B (KCC 2026)
     \item Sampling time-point introduced as a latent variable for temporal batch effects
           in tensile-strength prediction for biodegradable fiber spinning. With GBR,
           $R^2$ rose from 0.724 to 0.811, RMSE fell 17.2\%, and accuracy within the
           industry tolerance rose from 87.1\% to 92.4\%.
           MOTIE/KIAT project P0022335 (UCWIT 2025)
     \item Methodology for building software-in-the-loop verification environments
           at lower cost (2026, ongoing)
   \end{cvitems}}

% [확인] 제품명 Ark for CDC는 회사 제품 페이지(iarkdata.com)에서 확인. 담당 제품이 맞는지 확인 필요.
% DBMS 목록은 본인이 직접 분석한 것만. 제품이 지원하는 전체 목록을 옮기면 안 됨.
\cventry{2022.09--2024.06}{Researcher \cvnote{(reverse engineering)}}
  {Arkdata}{Daegu, Korea}
  {\begin{cvitems}
     \item Reverse-engineered transaction log formats across heterogeneous DBMSs
           (Oracle, PostgreSQL, MariaDB/MySQL, Tibero) to support development of a
           Change Data Capture (CDC) solution
     \item Designed and built a Kafka-based real-time pipeline across heterogeneous
           databases
     \item Implementation, analysis, and testing on a three-person research team
   \end{cvitems}}

% ============================================================ Teaching
% 6개 학기 7건. 프로그래밍 기초를 네 번 맡은 게 재위촉 신호라 굳이 설명 안 해도 읽힘.
\section{Teaching}
\cventry{2025.03--present}{Teaching Assistant}
  {School of Computer Science and Engineering, Kyungpook National University}{Daegu, Korea}
  {\begin{cvitems}
     \item Introduction to Programming: Spring 2025, Summer 2025, Winter 2025--26,
           Summer 2026 (four terms)
     \item System Programming and Creative Convergence Design: Fall 2025
     \item Java Programming: Spring 2026
   \end{cvitems}}

% ============================================================ Projects
\section{Projects}
% ABB는 기업이 아니라 경북대 ICT-ABB 융합전공. A(AI)-B(BigData)-B(Blockchain).
% 지역혁신사업 전자정보 DNA 사업단 소속 교육과정. 저장소 README에 "ABB 융합전공" 명시됨.
% 처음에 산학협력 과제로 잘못 적었던 것을 정정함 (ictabb.knu.ac.kr 확인).
% SIL 방법론 연구는 여기서 뺐음. 융합전공이 아니라 연구실 연구라 경력 섹션으로 옮김.
\cventry{2025--present}{ICT-ABB Convergence Major Project}
  {Kyungpook National University}
  {\cvlink{https://github.com/99JIK/ABB}{github.com/99JIK/ABB}}
  {Real-time emotion classification from facial action units. Up to 30 faces at once,
   IoU-based multi-face tracking, event triggers on emotion changes and AU thresholds,
   FastAPI REST and WebSocket streaming. MediaPipe, TensorFlow.}

% 기간은 저장소 커밋 활동(2026.01 demo - 2026.03)에서 추정. 실제 기간 다르면 수정.
\cventry{2026.01--2026.03}{ST Youngwon Smart Office}{Contract work, in production use}
  {\cvlink{https://github.com/99JIK/ST_YOUNGWON}{github.com/99JIK/ST\_YOUNGWON}}
  {\begin{cvitems}
     \item RAG question answering over internal documents. Regulatory documents chunked
           hierarchically (article, paragraph, fixed size), retrieved by hybrid keyword
           and vector search
     \item LLM providers (OpenAI, Claude, Gemini, local Ollama) behind a protocol
           abstraction so they can be swapped, per the client's maintainability requirement
     \item FastAPI, ChromaDB, SQLite, Synology NAS (FileStation API) integration,
           JWT auth, Docker Compose deployment
   \end{cvitems}}

% 수업 팀 프로젝트지만 내용이 실해서 넣음. ABB 2025년 표정 분석 주제와 결이 같아 서사도 이어짐.
% 기간은 과목코드(2025-2학기)와 저장소 최종 커밋(2025.12)에서 추정.
\cventry{2025.09--2025.12}{FOCUS \cvnote{(Facial Observation \& Concentration Understanding System)}}
  {Team coursework project}
  {\cvlink{https://github.com/99JIK/FOCUS_KNU_2025_COMP0738}{github.com/99JIK/FOCUS\_KNU\_2025\_COMP0738}}
  {Measures a learner's concentration from webcam video in real time. Scores distraction
   from five signals: eye closure, head orientation, gaze with head-rotation compensation,
   yawning, and facial action units. MediaPipe Face Landmarker, vanilla JS, Vite.}

% 제작 시기는 저장소 최종 커밋(2025.06)에서 추정. 2025-1학기 기말 무렵.
\cventry{2025}{Code Plagiarism Detector}{Personal project, Python}
  {\cvlink{https://github.com/99JIK/COPY_JIKILLER}{github.com/99JIK/COPY\_JIKILLER}}
  {Compares abstract syntax trees to judge similarity between submissions, ignoring
   variable names and comments. Supports Python (built-in \texttt{ast}), C/C++ (clang),
   and Java (javalang). Adjustable similarity threshold, side-by-side diff viewer.
   Built for the programming exams I assisted.}

% TIL 아카이브는 프로젝트에서 뺐음. 자동화 없는 블로그라 프로젝트로는 약하고,
% 링크는 이미 머리말에 있음.

% ============================================================ Technical Skills
\section{Technical Skills}
% [확인] GitHub 프로필의 스택 목록 그대로. 실무에서 안 쓴 건 빼는 게 낫다
% 국문판과 행 구성을 맞췄음. ML 항목은 UCWIT 2025에서 실제 사용한 근거 있음.
\cvline{Languages}{Python, Java, C/C++, TypeScript, JavaScript}
\cvline{Web / Backend}{FastAPI, Spring, Node.js, React, Vue.js}
\cvline{Data / ML}{Oracle, PostgreSQL, MySQL, ChromaDB, scikit-learn, LightGBM}
\cvline{Tools}{Docker, Git, Linux, Kafka}

% 수상 섹션은 뺐음. 2012 정보올림피아드 지역예선 동상뿐인데 시점이 14년 전이고 증거가 없어
% 지금 CV에 도움이 안 됨. 나중에 채울 게 생기면 되살릴 것.

\end{document}
